\begin{table}[p]\centering
\caption{Sources of the gains from levered diversification}\label{tab:composition-value}
\begin{minipage}{\linewidth}\small The table reports the change in the equivalent savings rate, in percentage points, from changing the asset class weights of a strategy at 175\% gross exposure. A negative change means that the household needs to save less to attain the expected utility of the all-equity strategy at 100\% exposure and a 10\% savings rate. Panel A adds global bonds, gold, managed futures, or all three to the levered all-equity strategy at their proportional-strategy weights, financed by reducing stocks. Panels B and C remove one asset class from the proportional or risk parity strategy and reallocate its weight to stocks. Reallocation is the weight moved, as a percentage of wealth. Panel D reports the weights that maximize expected utility at 175\% exposure in each sample, measured relative to the all-equity strategy at 175\%; the last two rows reoptimize after subtracting 3\% and 6\% per year from managed-futures returns. Each comparison is based on 10,000 bootstrap simulations. The last column draws returns from 1970 to 2025, after the 1968 separation of the free London gold market from the official price. In the full sample, the optimal weights are 83\% stocks, 92\% managed futures; 113\% stocks, 5\% gold, 56\% managed futures with a 3\% haircut; 127\% stocks, 2\% global bonds, 35\% gold, 11\% managed futures with a 6\% haircut. In the post-1970 sample, they are 72\% stocks, 18\% gold, 85\% managed futures; 96\% stocks, 30\% gold, 49\% managed futures with a 3\% haircut; 110\% stocks, 4\% global bonds, 43\% gold, 17\% managed futures with a 6\% haircut.\end{minipage}
\medskip

{\setlength{\tabcolsep}{4pt}\small
\begin{tabular}{lrrr}\toprule
Change in weights & Reallocation & \shortstack{Full sample\\ (1927--2025)} & \shortstack{Post-1970 sample\\ (1970--2025)} \\ \midrule
\multicolumn{4}{c}{Panel A: Add one asset class to the all-equity strategy at 175\%} \\
$+$ Global bonds & 46.7 & $-$2.1 & $-$4.7 \\
$+$ Gold & 29.2 & $-$3.3 & $-$7.3 \\
$+$ Managed futures & 29.2 & $-$5.4 & $-$8.3 \\
$+$ All three asset classes & 105.0 & $-$3.8 & $-$8.3 \\
\midrule
\multicolumn{4}{c}{Panel B: Remove one asset class from the proportional strategy at 175\%} \\
$-$ Global bonds & 46.7 & $-$1.7 & $-$1.0 \\
$-$ Gold & 29.2 & $-$0.8 & +0.4 \\
$-$ Managed futures & 29.2 & +1.2 & +1.1 \\
$-$ Bonds and gold & 75.8 & $-$1.6 & +0.0 \\
$-$ All three asset classes & 105.0 & +3.8 & +8.3 \\
\midrule
\multicolumn{4}{c}{Panel C: Remove one asset class from the risk parity strategy at 175\%} \\
$-$ Global bonds & 83.1 & $-$4.7 & $-$3.0 \\
$-$ Gold & 21.1 & $-$1.4 & $-$0.1 \\
$-$ Managed futures & 45.4 & +0.9 & +1.1 \\
$-$ Bonds and gold & 104.2 & $-$4.9 & $-$2.6 \\
$-$ All three asset classes & 149.6 & +1.4 & +6.8 \\
\midrule
\multicolumn{4}{c}{Panel D: Utility-maximizing weights at 175\% exposure} \\
Optimal mix & 92.2 & $-$6.9 & $-$10.2 \\
Optimal mix, MF $-$300bp & 61.7 & $-$4.6 & $-$8.6 \\
Optimal mix, MF $-$600bp & 47.6 & $-$3.5 & $-$8.0 \\
\bottomrule\end{tabular}}
\end{table}
